% ============================================================
%  MY UNDERSTANDING NOTES
%  Topic: Indian Children Speech Recognition
%  Base Paper: CARE (Dutta & Ganapathy, IEEE TASLP 2025)
% ============================================================

\documentclass[12pt, a4paper]{article}

\usepackage[margin=2.5cm]{geometry}
\usepackage{amsmath, amssymb}
\usepackage{xcolor}
\usepackage{tikz}
\usepackage{booktabs}
\usepackage{array}
\usepackage{titlesec}
\usepackage{mdframed}
\usepackage{enumitem}
\usepackage[hidelinks]{hyperref}

\usetikzlibrary{shapes.geometric, arrows.meta, positioning,
                fit, backgrounds, calc}

% ── Colors ──────────────────────────────────────────────────
\definecolor{pathA}{RGB}{41, 128, 185}
\definecolor{pathB}{RGB}{192, 57, 43}
\definecolor{pathC}{RGB}{39, 174, 96}
\definecolor{frozen}{RGB}{200, 200, 200}
\definecolor{trains}{RGB}{241, 196, 15}
\definecolor{lossclr}{RGB}{230, 126, 34}
\definecolor{notebg}{RGB}{255, 249, 230}
\definecolor{headclr}{RGB}{44, 62, 80}

% ── Section styling ─────────────────────────────────────────
\titleformat{\section}{\large\bfseries\color{headclr}}
  {\thesection}{1em}{}[\titlerule]
\titleformat{\subsection}{\normalsize\bfseries\color{headclr}}
  {\thesubsection}{1em}{}

% ── Note box ────────────────────────────────────────────────
\newmdenv[backgroundcolor=notebg, linecolor=orange!60,
          linewidth=1.5pt, roundcorner=5pt,
          innertopmargin=6pt, innerbottommargin=6pt]{notebox}

% ── TikZ styles ─────────────────────────────────────────────
\tikzset{
  box/.style={rectangle, rounded corners=4pt,
    minimum width=3.2cm, minimum height=0.65cm,
    text centered, draw, thick, font=\small},
  frozen/.style={box, fill=frozen!60, draw=gray},
  trains/.style={box, fill=trains!80, draw=orange!80},
  outbox/.style={box, font=\small\bfseries},
  lossbox/.style={box, fill=lossclr!50, draw=lossclr!90,
    font=\small\bfseries},
  arr/.style={->, >=Stealth, thick, line width=1.2pt}
}

% ── Title ───────────────────────────────────────────────────
\title{
  \textbf{My Research Notes}\\[4pt]
  \large Indian Children Speech Recognition\\[2pt]
  \normalsize Base Paper: CARE --- Dutta \& Ganapathy, IEEE TASLP 2025
}
\author{[Your Name]}
\date{\today}

% ============================================================
\begin{document}
\maketitle
\tableofcontents
\newpage

% ============================================================
\section{About the Project}
% ============================================================

\begin{itemize}[leftmargin=*]
  \item \textbf{Goal:} Build an ASR system specifically for Indian
        children's speech in Hindi, Marathi, and Indian English.
  \item \textbf{Dataset:} ASER --- 123.72 hours, 81,423 clips,
        5,260 children aged 5--18.
  \item \textbf{Base Paper:} CARE framework (SER) adapted for ASR.
  \item \textbf{Target:} ICASSP 2026.
\end{itemize}

% ============================================================
\section{ASER Dataset Summary}
% ============================================================

\begin{center}
\begin{tabular}{ll}
\toprule
\textbf{Property} & \textbf{Value} \\
\midrule
Total clips        & 81,423 \\
Unique children    & 5,260 \\
Total duration     & 123.72 hours \\
Avg.\ clip length  & 5.47 sec \\
Age range          & 5--18 years \\
\midrule
Hindi              & 15,880 clips — 53.84 hrs \\
English (Indian)   & 60,867 clips — 47.19 hrs \\
Marathi            &  4,421 clips — 22.39 hrs \\
\midrule
Rajasthan (Hindi)  & 31,114 clips — 46.17 hrs \\
UP (Hindi)         & 31,509 clips — 44.09 hrs \\
Maharashtra        & 18,800 clips — 33.45 hrs \\
\midrule
Correct readings   & 65,008 (79.8\%) \\
Incorrect readings & 16,404 (20.2\%) \\
\midrule
ASR train split    & 10,939 clips — 58.16 hrs \\
ASR dev split      &  1,404 clips —  7.68 hrs \\
ASR test split     &  1,530 clips —  7.35 hrs \\
\bottomrule
\end{tabular}
\end{center}

\begin{notebox}
\textbf{Important:} Splits are speaker-independent (split by child\_id).
No child's voice appears in more than one split. This prevents speaker
leakage and ensures honest WER evaluation.
\end{notebox}

% ============================================================
\section{Understanding the CARE Paper}
% ============================================================

% ────────────────────────────────────────────────────────────
\subsection{What CARE Does (Simple Version)}
% ────────────────────────────────────────────────────────────

CARE has \textbf{two expert listeners} looking at the same audio:
\begin{itemize}[leftmargin=*]
  \item \textbf{Semantic expert} --- understands the \textit{meaning}
        of what is said (like a literature professor reading the words)
  \item \textbf{Acoustic expert} --- understands \textit{how} it is
        said (pitch, speed, energy --- like a music producer)
\end{itemize}

Both experts are trained at the same time. Their outputs are merged
for the final task (emotion detection in CARE; text output in our case).

% ────────────────────────────────────────────────────────────
\subsection{Phase 1: Pre-Training}
% ────────────────────────────────────────────────────────────

\textbf{Paper location:} Section III-B, pages 3680--3681.

\medskip
\textbf{Running example:} Audio clip --- a child saying
\textit{``Radha ke paas ek tota hai''} (5 seconds, 80,000 numbers).

\medskip
\textbf{Key idea:} Three things happen to the same audio
simultaneously. Two paths produce the \textit{correct answers}
(teachers). One path produces the \textit{model's attempt} (student).

\bigskip

% ── THE BIG PHASE 1 DIAGRAM ─────────────────────────────────
\begin{center}
\begin{tikzpicture}[node distance=0.55cm and 2.2cm]

% ── Shared input ────────────────────────────────────────────
\node[box, fill=gray!20, draw=gray!70, minimum width=10cm]
  (audio) {\textbf{RAW AUDIO} \quad
           $\mathbf{x} \in \mathbb{R}^{80000}$
           \quad (same clip goes into all 3 paths)};

% ══════════ PATH A (left) ═══════════════════════════════════
\node[frozen, below left=1.2cm and 3.6cm of audio]
  (whisper) {Whisper ASR\\{\scriptsize frozen}};

\node[box, fill=pathA!15, draw=pathA,
      below=0.5cm of whisper]
  (text) {``Radha ke paas ek tota hai''\\
          {\scriptsize plain text string}};

\node[frozen, below=0.5cm of text]
  (roberta) {RoBERTa (full, frozen)\\
             {\scriptsize text model}};

\node[box, fill=pathA!15, draw=pathA,
      below=0.5cm of roberta]
  (r12) {12 rows $\times$ 768 numbers\\
         {\scriptsize one row per layer}};

\node[frozen, below=0.5cm of r12]
  (avgA) {Average Pool\\
          {\scriptsize 12 rows $\to$ 1 row}};

\node[outbox, fill=pathA!50, draw=pathA,
      below=0.5cm of avgA]
  (ytext) {$\mathbf{y}_\text{text} \in \mathbb{R}^{768}$\\
           {\scriptsize \textbf{CORRECT SEMANTIC ANSWER}}};

\draw[arr, color=pathA] (whisper) -- (text);
\draw[arr, color=pathA] (text)    -- (roberta);
\draw[arr, color=pathA] (roberta) -- (r12);
\draw[arr, color=pathA] (r12)     -- (avgA);
\draw[arr, color=pathA] (avgA)    -- (ytext);

% Path A label
\node[above=0.1cm of whisper, font=\small\bfseries,
      color=pathA] {PATH A --- Semantic Teacher};

% ══════════ PATH B (right) ══════════════════════════════════
\node[frozen, below right=1.2cm and 3.6cm of audio]
  (pase) {PASE+ (frozen)\\
          {\scriptsize voice quality extractor}};

\node[box, fill=pathB!15, draw=pathB,
      below=0.5cm of pase]
  (f50) {50 rows $\times$ 256 numbers\\
         {\scriptsize 1 row every 20ms}};

\node[frozen, below=0.5cm of f50]
  (down) {Downsample $\div 2$\\
          {\scriptsize keep every other row}};

\node[outbox, fill=pathB!50, draw=pathB,
      below=0.5cm of down]
  (ypase) {$\mathbf{y}_\text{pase} \in \mathbb{R}^{25 \times 256}$\\
           {\scriptsize \textbf{CORRECT ACOUSTIC ANSWER}}};

\draw[arr, color=pathB] (pase) -- (f50);
\draw[arr, color=pathB] (f50)  -- (down);
\draw[arr, color=pathB] (down) -- (ypase);

\node[above=0.1cm of pase, font=\small\bfseries,
      color=pathB] {PATH B --- Acoustic Teacher};

% ══════════ PATH C (centre) ═════════════════════════════════
\node[trains, below=1.2cm of audio]
  (cnn) {CNN Feature Extractor\\
         {\scriptsize WavLM conv layers --- TRAINS}};

\node[box, fill=pathC!15, draw=pathC,
      below=0.5cm of cnn]
  (cnnout) {$\mathbb{R}^{25 \times 768}$\\
            {\scriptsize 25 frames, 768 dim each}};

\node[trains, below=0.5cm of cnnout]
  (common) {Common Encoder\\
            {\scriptsize WavLM layers 1--6 --- TRAINS}\\
            {\scriptsize phonemes, pitch, basic sounds}};

\node[box, fill=pathC!15, draw=pathC,
      below=0.5cm of common]
  (comout) {$\mathbb{R}^{25 \times 768}$\\
            {\scriptsize richer representations}};

% Branch split
\node[trains, fill=pathB!40, draw=pathB,
      below left=0.8cm and 1.5cm of comout]
  (acbr) {Acoustic Branch\\
          {\scriptsize WavLM layers 7--12 + FC}\\
          {\scriptsize TRAINS}};

\node[trains, fill=pathA!40, draw=pathA,
      below right=0.8cm and 1.5cm of comout]
  (sembr) {Semantic Branch\\
           {\scriptsize RoBERTa layers}\\
           {\scriptsize + Conv Adapters --- TRAINS}};

\node[outbox, fill=pathB!40, draw=pathB,
      below=0.6cm of acbr]
  (yachat) {$\hat{\mathbf{y}}_\text{acoust}
             \in \mathbb{R}^{25 \times 256}$\\
            {\scriptsize Model's acoustic attempt}};

\node[outbox, fill=pathA!40, draw=pathA,
      below=0.6cm of sembr]
  (ysemhat) {$\hat{\mathbf{y}}_\text{sem}
              \in \mathbb{R}^{768}$\\
             {\scriptsize Model's semantic attempt}};

\draw[arr, color=pathC] (cnn)    -- (cnnout);
\draw[arr, color=pathC] (cnnout) -- (common);
\draw[arr, color=pathC] (common) -- (comout);
\draw[arr, color=pathB] (comout) -- (acbr);
\draw[arr, color=pathA] (comout) -- (sembr);
\draw[arr, color=pathB] (acbr)   -- (yachat);
\draw[arr, color=pathA] (sembr)  -- (ysemhat);

\node[above=0.1cm of cnn, font=\small\bfseries,
      color=pathC] {PATH C --- CARE Model (Student)};

% ══════════ Connect audio to all 3 paths ════════════════════
\draw[arr, color=gray!70]
  (audio.south) -- ++(0,-0.4) -| (whisper.north);
\draw[arr, color=gray!70]
  (audio.south) -- ++(0,-0.4) -- (cnn.north);
\draw[arr, color=gray!70]
  (audio.south) -- ++(0,-0.4) -| (pase.north);

% ══════════ LOSS CALCULATION ════════════════════════════════
\node[lossbox,
      below=0.7cm of $(ytext)!0.5!(yachat)$,
      xshift=1cm, minimum width=4.5cm]
  (lsem) {$\mathcal{L}_\text{sem}$ = MSE$(
           \mathbf{y}_\text{text},\,
           \hat{\mathbf{y}}_\text{sem})$\\
          {\scriptsize one number}};

\node[lossbox,
      below=0.7cm of $(ypase)!0.5!(ysemhat)$,
      xshift=-1cm, minimum width=4.5cm]
  (lacust) {$\mathcal{L}_\text{acoust}$ = MSE$(
             \mathbf{y}_\text{pase},\,
             \hat{\mathbf{y}}_\text{acoust})$\\
            {\scriptsize one number}};

\node[lossbox, fill=lossclr!80, draw=red!80,
      below=0.7cm of $(lsem)!0.5!(lacust)$,
      minimum width=6cm]
  (ltot) {$\mathcal{L}_\text{total} =
           \mathcal{L}_\text{sem} +
           \mathcal{L}_\text{acoust}}$\\[2pt]
         {\scriptsize Update PATH C only.
          Teachers (A, B) stay frozen.}};

\draw[arr, color=pathA]   (ytext)   |- (lsem);
\draw[arr, color=pathA]   (ysemhat) |- (lsem);
\draw[arr, color=pathB]   (ypase)   |- (lacust);
\draw[arr, color=pathB]   (yachat)  |- (lacust);
\draw[arr, color=lossclr] (lsem)    -- (ltot);
\draw[arr, color=lossclr] (lacust)  -- (ltot);

% ══════════ Legend ══════════════════════════════════════════
\node[below=0.7cm of ltot, font=\small, align=left]
  (legend) {
    \textcolor{gray}{$\blacksquare$} Gray box = Frozen (never updates) \quad
    \textcolor{orange}{$\blacksquare$} Yellow box = Trains (updates weights)\\[2pt]
    \textcolor{pathA}{$\longrightarrow$} Blue = Semantic path \quad
    \textcolor{pathB}{$\longrightarrow$} Red = Acoustic path \quad
    \textcolor{pathC}{$\longrightarrow$} Green = CARE model path
  };

\end{tikzpicture}
\end{center}

\bigskip

% ── Output summary table ─────────────────────────────────────
\begin{center}
\begin{tabular}{llll}
\toprule
\textbf{Symbol} & \textbf{Shape} & \textbf{From} &
\textbf{What it means} \\
\midrule
$\mathbf{y}_\text{text}$ & 768 numbers (1 row) &
  RoBERTa \textit{frozen} &
  Correct semantic answer --- sentence meaning \\[3pt]
$\mathbf{y}_\text{pase}$ & 25 rows $\times$ 256 &
  PASE+ \textit{frozen} &
  Correct acoustic answer --- voice quality over time \\[3pt]
$\hat{\mathbf{y}}_\text{sem}$ & 768 numbers (1 row) &
  Semantic branch \textit{trains} &
  Model's attempt at meaning \\[3pt]
$\hat{\mathbf{y}}_\text{acoust}$ & 25 rows $\times$ 256 &
  Acoustic branch \textit{trains} &
  Model's attempt at voice quality \\
\bottomrule
\end{tabular}
\end{center}

\begin{notebox}
\textbf{Phase 1 conclusion:} After training on thousands of clips,
$\mathcal{L}_\text{total} \to 0$.
CARE now produces both a 768-dim meaning vector AND a
$25\times256$ voice-quality grid from raw audio alone.
Teachers (RoBERTa and PASE+) are no longer needed.
\end{notebox}

% ────────────────────────────────────────────────────────────
\subsection{Phase 2: Fine-Tuning (Downstream Task)}
% ────────────────────────────────────────────────────────────

\textit{[To be added --- understanding in progress.]}

% ============================================================
\section{Our Proposed Architecture (Indian Children ASR)}
% ============================================================

\textit{[To be designed and added after full CARE understanding.]}

% ============================================================
\section{Planning and Next Steps}
% ============================================================

\textit{[To be added.]}

% ============================================================
\end{document}
